\documentclass[../DoAn.tex]{subfiles}
\begin{document}
\label{ch:ket_luan}
% =============================================================
% KẾT LUẬN (chương không đánh số)      Mục tiêu độ dài: ~2 trang
% =============================================================

\section*{Kết quả đạt được}              % ~0.8 trang

Đồ án đã đề xuất và hiện thực hoá \textit{CARL-VNE}, một phương pháp nhúng mạng ảo đa miền kết hợp mạng nơ-ron đồ thị với học tăng cường. Ba đóng góp chính đã được hoàn thành:

\begin{enumerate}
\item \textbf{Kiến trúc mã hoá phân cấp đa miền} gồm GCN nội miền và GCN liên miền, kết hợp đầu chọn ứng viên cross-attention tích hợp đặc trưng chi phí tường minh (cost-aware) và mặt nạ khả thi cứng, với toàn bộ mô hình chỉ khoảng $36{,}300$ tham số.
\item \textbf{Quy trình huấn luyện hai giai đoạn} (khởi động bằng học bắt chước từ heuristic MP-VNE rồi tinh chỉnh actor-critic có neo KL) cùng phát hiện then chốt rằng \textit{lựa chọn ứng viên} mới là đòn bẩy hiệu quả, và việc làm cho lựa chọn này ngẫu nhiên trong huấn luyện là điều kiện cần để vượt trần của mô hình bắt chước.
\item \textbf{Bộ benchmark đa bộ dữ liệu và đánh giá hệ thống} trên bộ dữ liệu $50$ nút với $21$ bộ dữ liệu phân phối VNR. Kết quả cho thấy phương pháp đề xuất vượt heuristic MP-VNE ở $19/21$ bộ dữ liệu phân phối (tại bộ dữ liệu trung tâm đạt $32{,}7\%$ so với $25{,}5\%$), cho thấy ưu thế bền vững trên một dải điều kiện vận hành rộng.
\end{enumerate}

\section*{Hạn chế}                      % ~0.6 trang

Phương pháp còn một số hạn chế. Thứ nhất, đánh giá hiện mới giới hạn ở quy mô $50$ nút; hiệu năng và khả năng mở rộng lên $100$, $200$ nút trở lên cần được kiểm chứng thêm. Thứ hai, ưu thế của phương pháp thu hẹp ở các vùng cực trị (tài nguyên khan hiếm, ràng buộc miền siết chặt, hoặc VNR rất lớn), nơi dư địa cải thiện nhỏ. Thứ ba, hiệu năng phụ thuộc phân phối VNR lúc huấn luyện; với phân phối lệch nhiều, chính sách có thể cần tinh chỉnh lại. Thứ tư, chi phí huấn luyện (khoảng $150$ epoch) tuy chấp nhận được nhưng không nhỏ, và đồ án chưa khảo sát đầy đủ khía cạnh công bằng tài nguyên giữa nhiều người thuê.

\section*{Hướng phát triển trong tương lai}  % ~0.6 trang

Một số hướng phát triển khả thi gồm: (1) \textbf{mở rộng quy mô} đánh giá lên $100$, $200$ nút rồi tới $500$–$1000$ nút (các tập dữ liệu đã sẵn sàng) với bộ mã hoá GAT đa lớp biểu cảm hơn; (2) \textbf{học chuyển giao và meta-RL} để chính sách thích nghi nhanh giữa các tô-pô và phân phối yêu cầu khác nhau; (3) \textbf{tích hợp cận trên ILP} cho các tô-pô nhỏ nhằm định lượng khoảng cách tới tối ưu; (4) \textbf{tối ưu trực tuyến} có tính tới thời gian sống và lịch đến của VNR để giữ ``khoảng trống chiến lược''; và (5) bổ sung \textbf{tiêu chí công bằng} giữa các miền và người thuê vào hàm thưởng. Việc hoàn tất so sánh trực tiếp giữa chế độ actor-critic và biến thể PPO có cắt cũng là một bước tiếp theo trực tiếp.

\end{document}
